\subsubsection{最少资源数与可执行资源链}
最少配置除了数量下界，还需要一份能实现它的实体安排。将某组某类任务按开始时刻排序，维护各已配置资源的最早释放时刻：能复用时选择一个已释放资源，否则新增同型实体。半开区间下，同刻释放先于开始处理。优先队列实现的时间复杂度为$O(n\log n)$。

若处理某任务时必须新增第$k$个实体，说明已有$k-1$个实体均未释放，加上当前任务，该时刻至少有$k$项并发占用。任意方案都需要至少$k$个，构造又恰好用到$k$个，所以最终数量等于最大重叠数。这与区间图的结构性质对应\cite{fulkerson1965interval}，也说明本题不必用通用启发式近似固定日程配置。

独立核验还可把满足前项释放不晚于后项开始的任务相连，形成可接续DAG。相应二部图最大匹配为$M$时，最小路径覆盖数为$|\mathcal J|-|M|$，每条路径对应一件实体资源。区间峰值、构造链和匹配分别从并发下界、可执行上界和图优化角度得到同一配置，形成三种计算结构的交叉支持。

对电池同样可输出“任务占用—恢复—下一任务”的具体链。因此资源配置不是抽象总量预测，而是具有明确任务归属的整数安排。数量、逐组向量和资源链同时保存，有助于部署阶段直接分配设备。
